% Cableado de la PCB_V2 según el esquema (redes de KiCad en letra pequeña)
\begin{tikzpicture}[
  font=\sffamily\small,
  bornera/.style={draw=grafito, fill=lima!35, rounded corners=1pt, align=center, minimum width=27mm, minimum height=10mm},
  pieza/.style={draw=grafito, fill=marfil, rounded corners=2pt, align=center, minimum height=10mm, inner sep=4pt},
  fusible/.style={draw=grafito, fill=white, align=center, minimum width=13mm, minimum height=6mm, font=\sffamily\scriptsize},
  aviso/.style={draw=rojoaviso, fill=rojoaviso!8, align=center, minimum width=13mm, minimum height=8mm, font=\sffamily\scriptsize},
  v12/.style={-{Stealth[length=2mm]}, thick, rojoaviso!80!black},
  dat/.style={-{Stealth[length=2mm]}, thick, azulnota},
  red/.style={font=\ttfamily\scriptsize, inner sep=1.5pt},
]
  % --- Entrada 1: regulador e inyección ---
  \node[bornera] (in1) at (0,0) {S\_LM\_INY1\\+12 V · GND};
  \node[fusible] (f5) at (3.6,0) {F5};
  \node[pieza] (lm) at (7.2,0) {LM2596\\12 V → \dato{reguladorSalida}};
  \node[pieza] (esp) at (12.2,0) {ESP32-S3\\5V / GND};
  \node[fusible] (f1) at (3.6,-1.6) {F1};
  \node[fusible] (f2) at (3.6,-2.8) {F2};
  \node[bornera] (iny1) at (7.2,-1.6) {INY1 · inyección 1};
  \node[bornera] (iny2) at (7.2,-2.8) {INY2 · inyección 2};
  \draw[v12] (in1) -- node[above, red, pos=0.45]{+12V\_INY} (f5);
  \draw[v12] (f5) -- node[above, red]{+12V\_LM} (lm);
  \draw[v12] (lm) -- node[above, red]{+5V\_ESP} (esp);
  \fill[rojoaviso!80!black] (2.25,0) circle (0.6mm);
  \draw[v12] (2.25,0) |- (f1.west);
  \draw[v12] (2.25,0) |- (f2.west);
  \draw[v12] (f1) -- (iny1);
  \draw[v12] (f2) -- (iny2);

  % --- Entrada 2: caras ---
  \node[bornera] (in2) at (0,-5.4) {S\_A\_B1\\+12 V · GND};
  \node[fusible] (f6) at (3.6,-4.8) {F6};
  \node[fusible] (f7) at (3.6,-6.0) {F7};
  \node[bornera] (a1) at (7.2,-4.8) {A1 → cara A\\GND · D1 · +12V};
  \node[bornera] (b1) at (7.2,-6.0) {B1 → cara B\\GND · D2 · +12V};
  \draw[v12] (in2.east) -- ++(4mm,0) |- (f6.west);
  \draw[v12] (in2.east) -- ++(4mm,0) |- (f7.west);
  \draw[v12] (f6) -- node[above, red]{+12V\_A} (a1);
  \draw[v12] (f7) -- node[below, red]{+12V\_B} (b1);

  % --- Datos ---
  \node[pieza, align=center] (b2) at (12.2,-5.4) {Bornera B2\\GPIO\_A · GPIO\_B · GPIO\_C\\R1 → D1 · R2 → D2};
  \draw[dat] (esp.south) -- node[right, font=\sffamily\scriptsize, align=left]{señales\\de datos} (b2.north);
  \draw[dat] (b2.west) ++(0,1.5mm) -- ++(-5mm,0) |- (a1.east);
  \draw[dat] (b2.west) ++(0,-1.5mm) -- ++(-5mm,0) |- (b1.east);
  \node[font=\sffamily\scriptsize, align=center] at (12.2,-6.85) {Poste que funciona: cara A en GPIO\dato{pinesPosteActual.caraA}, cara B en GPIO\dato{pinesPosteActual.caraB}\\ESP32 soldada en la \dato{pcbVersion}: GPIO\dato{pinesEsquema.GPIO_A} y GPIO\dato{pinesEsquema.GPIO_B}};
\end{tikzpicture}
